\documentclass[11pt,a4paper]{article}
\usepackage[utf8]{inputenc}
\usepackage[spanish]{babel}
\usepackage[margin=2.5cm]{geometry}
\usepackage{graphicx}
\usepackage{hyperref}
\usepackage{listings}
\usepackage{xcolor}
\usepackage{longtable}
\usepackage{booktabs}
\usepackage{titlesec}
 
\definecolor{codebg}{RGB}{245,245,245}
\lstset{
  backgroundcolor=\color{codebg},
  basicstyle=\ttfamily\small,
  breaklines=true,
  frame=single,
  columns=fullflexible
}
 
\title{Sistema de Comercio Electrónico con Arquitectura Hexagonal \\ \large Usuario, Producto y Pedido -- Backend Node.js + PostgreSQL, Frontend React}
\author{}
\date{\today}
 
\begin{document}
 
\maketitle
\tableofcontents
\newpage
 
\section{Objetivo General}
 
Construir la arquitectura conceptual y técnica de un sistema de comercio electrónico, modelando la base de datos y desarrollando operaciones CRUD completas para tres entidades interconectadas (Usuario, Producto, Pedido), con énfasis en la seguridad de credenciales y la separación estricta de capas mediante arquitectura hexagonal.
 
\section{Diagrama Arquitectónico}
 
El backend está organizado en cuatro capas siguiendo el patrón \textbf{Hexagonal (Ports \& Adapters)}, aplicado de forma repetida a las tres entidades del dominio, más un módulo transversal de Captcha.
 
\subsection{Capas del sistema}
 
\begin{longtable}{p{3cm}p{6cm}p{5cm}}
\toprule
\textbf{Capa} & \textbf{Responsabilidad} & \textbf{No conoce} \\
\midrule
\endhead
Domain & Entidades y reglas de negocio puras (validaciones, estados) & HTTP, SQL, Express, PostgreSQL \\
Application & Casos de uso (registrar, aprobar, listar) que orquestan el dominio & Detalles de PostgreSQL o Express \\
Infrastructure & Adaptadores de salida: implementan los puertos con PostgreSQL real & Reglas de negocio \\
Interfaces & Adaptadores de entrada: rutas Express que traducen HTTP en llamadas a casos de uso & Cómo se guardan los datos \\
\bottomrule
\end{longtable}
 
\subsection{Flujo de dependencias}
 
\begin{verbatim}
Cliente (React)
     |
     v
Interfaces (Controllers: user, producto, pedido, captcha)
     |
     v
Application (Services: UserService, ProductoService,
             PedidoService, CaptchaService)
     |
     v
Domain (Entidades: User, Producto, Pedido, Captcha
        + Puertos: *RepositoryPort)
     ^
     | implementa
     |
Infrastructure (Adapters: *RepositoryAdapter -> PostgreSQL
                CaptchaStoreAdapter -> memoria)
\end{verbatim}
 
La inyección de dependencias ocurre en \texttt{server.js} (composition root): ahí se instancian los adaptadores concretos y se inyectan en los servicios de aplicación, que a su vez se inyectan en los controladores HTTP.
 
\section{Modelado de la Base de Datos (PostgreSQL)}
 
\subsection{Tabla usuarios}
\begin{lstlisting}[language=SQL]
CREATE TABLE usuarios (
  id SERIAL PRIMARY KEY,
  nombre VARCHAR(100) NOT NULL,
  email VARCHAR(150) NOT NULL UNIQUE,
  password VARCHAR(255) NOT NULL,
  rol VARCHAR(20),
  estado VARCHAR(20) NOT NULL DEFAULT 'pendiente',
  creado_en TIMESTAMP DEFAULT CURRENT_TIMESTAMP
);
\end{lstlisting}
 
\subsection{Tabla productos}
\begin{lstlisting}[language=SQL]
CREATE TABLE productos (
  id SERIAL PRIMARY KEY,
  nombre VARCHAR(150) NOT NULL,
  descripcion TEXT,
  precio NUMERIC(10,2) NOT NULL,
  imagen_url VARCHAR(255),
  vendedor_id INT REFERENCES usuarios(id),
  estado VARCHAR(20) NOT NULL DEFAULT 'pendiente',
  creado_en TIMESTAMP DEFAULT CURRENT_TIMESTAMP
);
\end{lstlisting}
 
\subsection{Tabla pedidos}
\begin{lstlisting}[language=SQL]
CREATE TABLE pedidos (
  id SERIAL PRIMARY KEY,
  producto_id INT REFERENCES productos(id),
  usuario_id INT REFERENCES usuarios(id),
  estado VARCHAR(20) NOT NULL DEFAULT 'pendiente',
  creado_en TIMESTAMP DEFAULT CURRENT_TIMESTAMP
);
\end{lstlisting}
 
\textbf{Nota sobre imágenes:} el campo \texttt{imagen\_url} almacena únicamente una referencia (URL) a la imagen del producto; el archivo binario nunca se guarda en la base de datos, cumpliendo con el requerimiento de la actividad.
 
\section{Especificación de Endpoints}
 
\subsection{Captcha}
\begin{longtable}{p{1.5cm}p{3cm}p{7cm}}
\toprule
\textbf{Método} & \textbf{Ruta} & \textbf{Descripción} \\
\midrule
\endhead
GET & /captcha & Genera un nuevo desafío de 10 imágenes SVG (una por carácter) \\
\bottomrule
\end{longtable}
 
\subsection{Usuario}
\begin{longtable}{p{1.5cm}p{4cm}p{7cm}}
\toprule
\textbf{Método} & \textbf{Ruta} & \textbf{Descripción} \\
\midrule
\endhead
POST & /usuarios/registro & Registra un usuario nuevo (queda \texttt{pendiente}, sin rol); requiere captcha \\
POST & /usuarios/login & Inicia sesión; requiere estar \texttt{aprobado} y tener rol asignado; requiere captcha \\
GET & /usuarios & Lista todos los usuarios (panel admin) \\
GET & /usuarios/pendientes & Lista usuarios en estado \texttt{pendiente} \\
PUT & /usuarios/:id/aprobar & El admin asigna rol y aprueba o rechaza \\
DELETE & /usuarios/:id & Elimina un usuario \\
\bottomrule
\end{longtable}
 
\subsection{Producto}
\begin{longtable}{p{1.5cm}p{4cm}p{7cm}}
\toprule
\textbf{Método} & \textbf{Ruta} & \textbf{Descripción} \\
\midrule
\endhead
POST & /productos & Registra un producto (rol \texttt{producto}); queda \texttt{pendiente} \\
GET & /productos & Lista todos los productos (panel admin) \\
GET & /productos/aprobados & Lista solo productos aprobados (vista Marketplace, rol \texttt{pedido}) \\
GET & /productos/pendientes & Lista productos pendientes de aprobación \\
PUT & /productos/:id & Actualiza datos del producto (dueño) \\
PUT & /productos/:id/aprobar & El admin aprueba o rechaza el producto \\
DELETE & /productos/:id & Elimina un producto \\
\bottomrule
\end{longtable}
 
\subsection{Pedido}
\begin{longtable}{p{1.5cm}p{4cm}p{7cm}}
\toprule
\textbf{Método} & \textbf{Ruta} & \textbf{Descripción} \\
\midrule
\endhead
POST & /pedidos & Crea un pedido (interés sobre un producto del catálogo) \\
GET & /pedidos & Lista todos los pedidos \\
PUT & /pedidos/:id & Actualiza estado (confirmar / cancelar) \\
DELETE & /pedidos/:id & Elimina un pedido \\
\bottomrule
\end{longtable}
 
\section{Estructura del Repositorio}
 
\subsection{Backend (market-service/)}
\begin{lstlisting}
market-service/
  src/
    domain/
      user.js, userRepositoryPort.js
      producto.js, productoRepositoryPort.js
      pedido.js, pedidoRepositoryPort.js
      captcha.js, captchaRepositoryPort.js
    application/
      userService.js, productoService.js
      pedidoService.js, captchaService.js
    infrastructure/
      db.js
      userRepositoryAdapter.js
      productoRepositoryAdapter.js
      pedidoRepositoryAdapter.js
      captchaStoreAdapter.js
    interfaces/
      userController.js, productoController.js
      pedidoController.js, captchaController.js
    server.js
  database.sql
  package.json
  .env
\end{lstlisting}
 
\subsection{Frontend (market-frontend/)}
\begin{lstlisting}
market-frontend/
  src/
    App.jsx                (login, registro, captcha, routing por rol)
    DashboardAdmin.jsx      (aprobar usuarios y productos)
    DashboardProducto.jsx   (CRUD de productos propios)
    DashboardPedido.jsx     (catalogo tipo marketplace)
    main.jsx
  vite.config.js
  package.json
\end{lstlisting}
 
\section{Guía de Despliegue en WSL}
 
\subsection{Prerrequisitos (una sola vez por máquina)}
\begin{lstlisting}[language=bash]
sudo apt update && sudo apt upgrade -y
 
# Node.js 20 LTS
curl -fsSL https://deb.nodesource.com/setup_20.x | sudo -E bash -
sudo apt install -y nodejs
 
# PostgreSQL
sudo apt install -y postgresql postgresql-contrib
 
# Git y utilidades
sudo apt install -y git unzip curl
\end{lstlisting}
 
\subsection{Cada vez que se abre WSL}
PostgreSQL no arranca automáticamente en WSL:
\begin{lstlisting}[language=bash]
sudo service postgresql start
\end{lstlisting}
 
\subsection{Configuración de la base de datos (una sola vez)}
\begin{lstlisting}[language=bash]
sudo -u postgres psql
\end{lstlisting}
\begin{lstlisting}[language=SQL]
CREATE USER apiuser WITH PASSWORD 'tu_password';
CREATE DATABASE market_db OWNER apiuser;
GRANT ALL PRIVILEGES ON DATABASE market_db TO apiuser;
\end{lstlisting}
\begin{lstlisting}[language=bash]
cd market-service
psql -U apiuser -d market_db -h localhost -f database.sql
\end{lstlisting}
 
\subsection{Levantar el Backend}
\begin{lstlisting}[language=bash]
cd market-service
npm install
# crear .env con DB_HOST, DB_PORT, DB_USER, DB_PASSWORD, DB_NAME, PORT
node src/server.js
\end{lstlisting}
Queda escuchando en \texttt{http://localhost:4000}.
 
\subsection{Levantar el Frontend}
En otra terminal:
\begin{lstlisting}[language=bash]
cd market-frontend
npm install
npm run dev
\end{lstlisting}
Queda escuchando en \texttt{http://localhost:5173}.
 
\subsection{Consumo en el navegador}
Como WSL comparte red con Windows, se accede directamente desde el navegador en:
\texttt{http://localhost:5173}. El frontend realiza peticiones \texttt{fetch} al backend en el puerto 4000, habilitado mediante CORS.
 
\section{Notas de Seguridad Implementadas}
 
\begin{itemize}
  \item \textbf{Contraseñas}: nunca se almacenan en texto plano; se cifran con \texttt{bcrypt} (10 salt rounds) antes de insertarse en PostgreSQL.
  \item \textbf{Captcha}: obligatorio en login y registro, de un solo uso, expira a los 5 minutos, almacenado en memoria (no en la base de datos) por tratarse de un dato temporal. Compuesto por 10 imágenes SVG generadas dinámicamente, una por carácter, con rotación y color aleatorios.
  \item \textbf{Aprobación de acceso}: ningún usuario puede iniciar sesión hasta que un administrador le asigna un rol (\texttt{admin}, \texttt{producto} o \texttt{pedido}) y aprueba su cuenta.
  \item \textbf{Aprobación de productos}: ningún producto es visible en el catálogo público hasta que un administrador lo aprueba.
  \item \textbf{Variables sensibles}: las credenciales de conexión a la base de datos se mantienen en un archivo \texttt{.env}, excluido del control de versiones mediante \texttt{.gitignore}.
\end{itemize}
 
\section{Conclusiones}
 
El sistema implementa exitosamente las tres entidades solicitadas (Usuario, Producto, Pedido) bajo arquitectura hexagonal, con separación estricta entre reglas de negocio (dominio), orquestación de casos de uso (aplicación), acceso a datos (infraestructura) y exposición HTTP (interfaces). El frontend en React consume estos endpoints y presenta una interfaz diferenciada según el rol del usuario autenticado, cumpliendo con el flujo de aprobación administrativa requerido por la actividad.
 
\end{document}
 